\documentclass[10pt,twocolumn]{article}

\usepackage[T1]{fontenc}
\usepackage[utf8]{inputenc}
\usepackage{lmodern}
\usepackage{microtype}
\usepackage{graphicx}
\usepackage{booktabs}
\usepackage{tabularx}
\usepackage{array}
\usepackage{amsmath}
\usepackage{amssymb}
\usepackage{siunitx}
\usepackage{subcaption}
\usepackage[colorlinks=true,linkcolor=blue,citecolor=blue,urlcolor=blue]{hyperref}
\captionsetup{hypcap=false}

\sisetup{
  round-mode=places,
  round-precision=3,
  detect-all
}

\title{\textbf{Adaptive Quality Map Hyperstacking for Pixel-Wise Quality-Weighted Deep-Sky Reconstruction}\\
\large Method definition v0.2.0, implementation contract, post-reconstruction robustness extensions, and an M31 validation run}
\author{Jeamy Lee$^{1}$ \\ \small $^{1}$Independent Researcher / tile\_compile project}
\date{July 22, 2026}

\begin{document}
\raggedbottom
\twocolumn[
\maketitle

\begin{center}
\begin{minipage}{0.88\textwidth}
\small
\textbf{Abstract}\par
\vspace{0.35em}
Adaptive Quality Map Hyperstacking (AQMH) is a deterministic pixel-wise quality-weighted
reconstruction method for short-exposure deep-sky imaging. It replaces block-constant local
quality scores with dense per-frame quality maps and reconstructs each supported output pixel by
a weighted mean whose AQMH weight is evaluated at that same pixel. The method is independent of
Classic Tile Compile: registration, photometric normalization, common-overlap handling, reporting,
and run management may be shared, but AQMH does not use classic local tile metrics, state
clustering, or synthetic frames as reconstruction inputs or fallbacks.
This paper consolidates the AQMH v0.2.0 definition, including the post-reconstruction robustness
extensions formalised in the v0.2.1 methodology: a background-gradient penalty in global frame
quality, a registration-weight guard, adaptive low-frequency neutralisation, and structure-masked
detail blending. It states the implementation contract used by tile\_compile and evaluates a
concrete M31 validation run. The run processed 645 frames on a $3924\times2310$ canvas, wrote 645
dense luminance quality maps at half resolution, and reconstructed the image with
\num{5.350e9} finite map samples, zero missing map samples, and no fallback to classic tile
weights. The measured median FWHM changed from \num{12.329} px to \num{5.268} px, corresponding
to a \num{57.273}\% reduction. The background RMS changed from \num{4.185} to \num{0.780},
corresponding to an \num{81.351}\% decrease. The uniform-control validation gate triggered on
the background-RMS regression (\num{0.168} $>$ \num{0.05} threshold); the raw AQMH output was
preserved by the guard. These measurements demonstrate a single-run execution of the AQMH path
with the full v0.2.0 extension suite active; they do not by themselves establish population-level
performance across object classes, instruments, or atmospheric regimes.
\end{minipage}
\end{center}

\vspace{0.35cm}
\noindent \textbf{Keywords:} astronomical image reconstruction, deep-sky imaging,
quality maps, pixel-wise weighted mean, short exposure stacking, deterministic reconstruction,
sigma clipping, uniform control validation
\vspace{0.6cm}
]

\section{Introduction}

Short-exposure deep-sky imaging collects many noisy frames instead of a small number of long
exposures. The data often contain useful signal even when individual frames suffer from local
defects: cloud edges, satellite traces, registration remnants, small hot-pixel clusters, or
spatially varying seeing. A single global frame score cannot describe these local variations.
Coarse tile scores improve locality, but they still assign one scalar quality value to an
extended region. This has two structural consequences. First, a small contaminated area can
depress the weight of a much larger block. Second, neighboring blocks can receive different
constant weights, creating a discontinuous weighting field that is unrelated to the optical
image content.

AQMH is documented here as an independent reconstruction path within the tile\_compile project.
The reference method for Classic Tile Compile is
\href{https://doi.org/10.5281/zenodo.19437794}{\emph{Tile-Based Quality Reconstruction for DSO},
Version 3.3.9}.

AQMH addresses the locality problem by computing a dense quality field
$Q_{f,c}^{\mathrm{map}}(x,y)$ for each frame $f$ and channel or stream $c$. The reconstruction
then uses the frame sample itself and the AQMH weight field at the same pixel:
\begin{equation}
  R_c(p) =
  \frac{\sum_f W_{f,c}^{\mathrm{aqmh}}(p)\, I_{f,c}(p)}
       {\sum_f W_{f,c}^{\mathrm{aqmh}}(p)} ,
  \quad W_{f,c}^{\mathrm{aqmh}}(p) \ge 0 .
\end{equation}
After the weights have been computed, the reconstruction core is therefore a deterministic linear
weighted mean. AQMH does not hallucinate pixel values, does not synthesize replacement frames, and
does not rely on tile clustering. Its purpose is narrower: retain useful local regions from every
registered frame while suppressing locally degraded regions through a dense, reproducible quality
field.

\section{Method overview}

\subsection{Inputs and independence}

AQMH reuses only shared preprocessing outputs:
\begin{itemize}
  \item calibrated, registered, prewarped, and normalized frames;
  \item a canvas-valid/common-overlap mask;
  \item global frame diagnostics used for a frame-level weight;
  \item run management, artifacts, reports, and UI status.
\end{itemize}
AQMH is otherwise independent from Classic Tile Compile. In particular, the run studied here
reports \texttt{classic\_tile\_weights\_used=false} and
\texttt{fallback\_to\_classic=false}. The classic phases \texttt{STATE\_CLUSTERING} and
\texttt{SYNTHETIC\_FRAMES} are not part of AQMH reconstruction.

The distinction matters operationally. Shared preprocessing produces registered, normalized image
samples and a valid canvas domain. AQMH then defines its own quality maps and reconstruction
weights. Any classic tile grid that appears in reports is therefore a geometry or diagnostic
artifact, not the AQMH weight lattice.

\subsection{Binding invariants}

The AQMH implementation follows eight practical invariants:
\begin{enumerate}
  \item No frame is discarded solely because of its AQMH quality map.
  \item Canvas-invalid pixels are excluded from all map and reconstruction statistics.
  \item Once weights are known, reconstruction is a linear weighted mean of finite samples.
  \item Quality maps and diagnostics are deterministic for identical input data and parameters.
  \item AQMH outputs weights, masks, and diagnostics only; it does not predict missing signal.
  \item Cherry-pick and any per-pixel selection mode must enforce a documented minimum
        retained-sample count, below which the mode is disabled automatically.
  \item Cross-infrastructure extensions (registration-weight guard, background penalty) may only
        multiply the global frame weight by a non-negative scalar. They must not alter
        $Q^{\mathrm{map}}$, must not introduce pixel-dependent weights, and must be independently
        disableable.
  \item Any post-reconstruction neutralisation or detail blending must be validation-gated. The
        run must record which candidate was selected and why. The final output must not be worse,
        by the documented regression metrics, than the ordinary AQMH output.
\end{enumerate}
The M31 run uses \texttt{cherry\_pick\_enabled=false}; therefore no optional pixel-level
frame subset mode is active.

\subsection{Quality-map model}

For a normalized registered input frame $I_{f,c}(x,y)$, AQMH computes a multi-scale quality
field. The default pyramid uses downscale factors $D_s=4^s$ with a local radius of four
pixels in scale space. Scales that are too coarse for the image dimensions are omitted and the
geometric-mean exponent is computed from the actually available scale count. At each retained
scale $s$, three local signals are evaluated over valid canvas support:
\begin{itemize}
  \item a sharpness signal based on local Laplacian-response variance;
  \item a local SNR signal based on background-centered positive signal over robust scale;
  \item an artifact anomaly gate based on high-pass outlier fraction.
\end{itemize}
The sharpness and SNR signals are normalized within a frame by robust z-scores and combined by a
sigmoid:
\begin{equation}
\begin{split}
  \Psi_s(x,y) &=
  \sigma\!\left(
    w_{\mathrm{sharp}}\, z(\Phi_{\mathrm{sharp},s})
    + w_{\mathrm{snr}}\, z(\Phi_{\mathrm{snr},s})
  \right) \\
  &\quad \times \Phi_{\mathrm{artifact},s}(x,y).
\end{split}
\end{equation}
Scale maps are upsampled with valid-support interpolation and fused by geometric mean:
\begin{equation}
  Q_{f,c}^{\mathrm{map}}(x,y) =
  \left(\prod_{s \in S_{\mathrm{actual}}} \Psi_s^{\mathrm{up}}(x,y)\right)^{1/|S_{\mathrm{actual}}|}.
\end{equation}
The geometric mean is intentionally conservative: one poor scale suppresses the fused quality at
that location, and a zero at any scale remains an exact veto after fusion. Canvas-invalid pixels
are set to zero only as an output convention after they have already been excluded from all
statistics.

The z-score normalization also fixes the interpretation of $Q^{\mathrm{map}}$. The sigmoid
component is a within-frame relative quality field; it indicates which regions of a frame are
locally better or worse than other regions of the same frame. Whole-frame discrimination is
carried by the global term $G_{f,c}$ and by absolute gates such as the artifact fraction.

\subsection{Global quality with background penalty}

The effective AQMH weight multiplies the frame-level quality term and the spatial quality map.
The v0.2.0 global quality formula introduces an optional background-gradient penalty as a third
signal:
\begin{equation}
\begin{split}
  \mathrm{score}_f &= w_{\mathrm{sharp}}\, z_{s,f}
    + w_{\mathrm{snr}}\, z_{n,f}
    - w_{\mathrm{bg}}\, z_{b,f} \\[2pt]
  G_{f,c} &= g_{\mathrm{floor}} + (1 - g_{\mathrm{floor}})\,
    \sigma\!\left(\mathrm{clamp}(k_{\mathrm{scale}} \cdot \mathrm{score}_f,\, -8,\, 8)\right)
\end{split}
\end{equation}
where $z_{b,f}$ is the robust z-score of the per-frame sky-gradient summary
$g_{\mathrm{background},f} = (q_{\max} - q_{\min}) / \mathrm{background}_f$, computed from
quadrant-based masked median backgrounds.

A strong large-scale gradient means that the background model used by the SNR signal is a biased
local estimate. Frames with very flat backgrounds tend to be cleaner and more linear; frames with
a strong moon-glow or light-pollution gradient receive a lower global weight, reducing their
influence on the final stack without removing them entirely.

The weights are renormalised so that $w_{\mathrm{sharp}} + w_{\mathrm{snr}} + w_{\mathrm{bg}} = 1$
when at least one effective weight is positive. If a configured weight is $\le 0$, if fewer than
three finite positive summary values are available, or if the coefficient of variation
$\mathrm{CV} = \mathrm{MAD}/\mathrm{median}$ is below $0.01$, the corresponding effective weight
is set to zero. This prevents near-constant signals from amplifying noise.

\subsection{Registration-weight guard}

Before reconstruction, the runner may optionally apply a registration-weight guard to $G_{f,c}$.
It reads the shared \texttt{global\_registration.json} artifact and damps the global weight of
frames with low cross-correlation confidence, sequential or predicted registration sources, or
deep chain-depth:
\begin{equation}
  G_{f,c}^{\mathrm{eff}} = G_{f,c} \cdot R_f, \quad R_f \in [r_{\mathrm{floor}},\, 1].
\end{equation}
The guard is a shared-infrastructure multiplier, not part of the AQMH quality model. It prevents
high AQMH quality scores from amplifying registration errors.

\subsection{Reconstruction}

For each output pixel, the implementation accumulates finite, canvas-valid samples only. Missing
quality maps are not replaced by unweighted mean contributions. A zero quality map is an
explicit veto for that frame/pixel contribution. Iterative sigma clipping with
$\sigma_{\mathrm{low}}$ and $\sigma_{\mathrm{high}}$ rejects outlier samples before the weighted
mean is computed. A minimum effective sample count $n_{\mathrm{eff}}^{\min}$ and minimum fraction
$f_{\min}$ are enforced.

The studied M31 run stores one luminance quality map per frame at half canvas resolution
(\texttt{resolution\_divisor=2}) and interpolates it for reconstruction. This is a storage and
I/O choice, not a change to the mathematical definition of AQMH as a dense pixel-weighted
reconstruction method.

\section{Post-reconstruction robustness extensions}

\subsection{Uniform control mean}

When \texttt{compute\_uniform\_control} is enabled, the reconstruction routine also computes the
unweighted mean of the registered frames over the same valid-pixel set:
\begin{equation}
  U(p) = \frac{1}{|V_c^I(p)|} \sum_{f \in V_c^I(p)} I_{f,c}(p).
\end{equation}
$U(p)$ is not a quality-weighted result; it is a diagnostic reference for the validation gates.

\subsection{Validation metrics}

Both the raw AQMH output $A(p)$ and any post-processed candidate are compared to $U(p)$ using:
\begin{itemize}
  \item FWHM regression: $\Delta_{\mathrm{FWHM}} = (\mathrm{FWHM}_{\mathrm{cand}} - \mathrm{FWHM}_{\mathrm{ctrl}}) / \mathrm{FWHM}_{\mathrm{ctrl}}$
  \item Background RMS regression: $\Delta_{\mathrm{RMS}} = (\mathrm{RMS}_{\mathrm{cand}} - \mathrm{RMS}_{\mathrm{ctrl}}) / \mathrm{RMS}_{\mathrm{ctrl}}$
  \item Seam score regression, tail$_{11}$ regression, elongation regression.
\end{itemize}
Each metric records an applicability state: \texttt{pass}, \texttt{fail}, or
\texttt{not\_applicable}. A metric is \texttt{not\_applicable} when either side of the comparison
is missing, non-finite, has too few samples, or would require a relative regression against a
zero or near-zero control denominator. A \texttt{not\_applicable} metric is diagnostic only and
must not by itself trigger fallback.

\subsection{Adaptive low-frequency neutralisation}

If $U(p)$ is available, an optional neutralised candidate $N(p)$ is computed:
\begin{equation}
\begin{split}
  L(p) &= \mathrm{GaussianBlur}(A(p) - U(p),\, \sigma = 96\,\mathrm{px}) \\
  N(p) &= A(p) - L(p).
\end{split}
\end{equation}
The blur uses reflect-101 border mode. The sigma of \SI{96}{px} targets large-scale ``veil''
residuals without affecting structure on smaller scales. $N(p)$ may be selected as the base over
$A(p)$ only when the background RMS regression is strictly smaller.

\subsection{Structure-masked detail blending}

After selecting the base $B(p)$ (either $A(p)$ or $N(p)$), a second optional candidate preserves
AQMH detail in high-structure regions while keeping the smoother background from $B(p)$:
\begin{equation}
  D(p) = U(p) + M_s(p) \cdot (B(p) - U(p))
\end{equation}
where $M_s(p)$ is a soft structure mask derived from the gradient magnitude of $U(p)$, mapped
from the $q_{\mathrm{low}} = 0.40$ quantile to $q_{\mathrm{high}} = 0.90$ quantile to $[0, 1]$
and Gaussian-blurred with $\sigma = \SI{4}{px}$. In high-structure regions ($M_s \approx 1$) the
candidate follows the AQMH base; in smooth regions ($M_s \approx 0$) it follows the uniform
control.

If $D(p)$ fails the applicable validation thresholds, an alpha-blended attenuation search is
performed over $\alpha \in [0, 1]$. The largest $\alpha$ that satisfies the thresholds is
selected. If no $\alpha > 0$ passes, the base $B(p)$ is kept.

\subsection{Final protection gate}

After all post-processing, the selected output $O(p)$ is validated one more time against $U(p)$
and the immutable raw AQMH baseline $A(p)$. If any applicable configured threshold is exceeded,
the runner rejects the post-processed candidate and emits $A(p)$. Blending toward $U(p)$ is
forbidden because it can erase faint diffuse signal while one-sided noise metrics appear to
improve.

\section{Implementation run}

\subsection{Data and configuration}

The evaluation data are taken from a tile\_compile AQMH run on an M31 dataset:
\begin{center}
\small
\path{runs/M31-default-afix1_20260722_111345}
\end{center}
The paper bundle itself is archived under
\href{https://github.com/jeamy/tile_compile/tree/master/docs/AQMH/zenodo-0.2.0}
{\path{docs/AQMH/zenodo-0.2.0}} in the project repository. Run artifacts used for
the measurements are local analysis outputs located in:
\begin{center}
\small
\path{runs/M31-default-afix1_20260722_111345/artifacts}
\end{center}
The final preview shown in Figure~\ref{fig:preview} is derived from
\path{outputs/stacked_rgb_hms.fits}.

This paper treats the run as an implementation validation case. It verifies that the expected AQMH
execution path is followed and reports the resulting diagnostics. It does not compare AQMH against
a separately tuned Classic Tile Compile run on the same data, and it does not claim that the
reported percentage changes generalize beyond this run.

\begin{table}[t]
\centering
\caption{M31 run geometry and frame counts.}
\begin{tabular}{lr}
\toprule
Quantity & Value \\
\midrule
Input frames loaded & 645 \\
Usable frames & 645 \\
Canvas size & $3924\times2310$ px \\
Common-overlap fraction & \num{0.9748} \\
Reconstruction fraction & \num{0.9823} \\
Adaptive tiles in shared grid & 209 \\
Tile size / stride & 296 px / 208 px \\
\bottomrule
\end{tabular}
\label{tab:geometry}
\end{table}

\begin{table}[t]
\centering
\caption{AQMH map storage and reconstruction support.}
{\small
\begin{tabular}{@{}lr@{}}
\toprule
Quantity & Value \\
\midrule
AQMH enabled & true \\
Frames/maps & 645 / 645 \\
Stored map size & $1962\times1155$ px \\
Resolution divisor & 2 \\
Map dtype & uint16 \\
Finite samples & \num{5.350e9} \\
Missing samples & 0 \\
Unsupported output pixels & \num{18539} \\
Zero-veto pixels & 0 \\
Classic fallback & false \\
Execution backend & cuda\_native\_v0\_2 \\
\bottomrule
\end{tabular}
}
\label{tab:support}
\end{table}

\begin{table}[t]
\centering
\caption{v0.2.0 extension configuration.}
{\small
\begin{tabular}{@{}lr@{}}
\toprule
Parameter & Value \\
\midrule
\multicolumn{2}{@{}l}{\emph{Global quality}} \\
$g_{\mathrm{floor}}$ & 0.03 \\
$w_{\mathrm{sharp}}$ & 0.55 \\
$w_{\mathrm{snr}}$ & 0.30 \\
$w_{\mathrm{bg}}$ & 0.25 \\
$k_{\mathrm{scale}}$ & 1.5 \\
\multicolumn{2}{@{}l}{\emph{Registration-weight guard}} \\
Enabled & true \\
$r_{\mathrm{floor}}$ & 0.30 \\
$\mathrm{cc}_{\mathrm{floor}}$ & 0.35 \\
$\mathrm{cc}_{\mathrm{full}}$ & 0.80 \\
Sequential factor & 0.92 \\
Predicted factor & 0.50 \\
\multicolumn{2}{@{}l}{\emph{Reconstruction}} \\
Clip iterations & 4 \\
$\sigma_{\mathrm{low}}$, $\sigma_{\mathrm{high}}$ & 2.0, 2.0 \\
$f_{\min}$ & 0.40 \\
$n_{\mathrm{eff}}^{\min}$ & 2.0 \\
\multicolumn{2}{@{}l}{\emph{Neutralisation}} \\
Blur $\sigma$ & \SI{96}{px} \\
\multicolumn{2}{@{}l}{\emph{Structure mask}} \\
$q_{\mathrm{low}}$, $q_{\mathrm{high}}$ & 0.40, 0.90 \\
Blur $\sigma$ & \SI{4}{px} \\
\bottomrule
\end{tabular}
}
\label{tab:config}
\end{table}

\subsection{Quality-map diagnostics}

The run wrote one luma quality map per frame. The average map mean was
\num{0.498315} with standard deviation \num{0.002051}. The per-frame map mean ranged from
\num{0.492467} to \num{0.503678}. The average p10 and p90 across frames were \num{0.376616}
and \num{0.635871} respectively. The average artifact fraction was \num{1.92e-6}, with a
range from \num{0} to \num{1.2e-5}. All 645 diagnostics reported the
\texttt{scene\_dependent\_snr} flag, which marks use of the degraded local-SNR path when the
background-centered support condition is not fully satisfied.

The global sharpness source was \texttt{psf\_wfwhm\_inverted} for all frames, meaning the
PSF-FWHM sharpness proxy was used instead of the Laplacian-variance signal. This avoids a known
inversion of \texttt{local\_variance(laplacian)} on extended-emission targets such as M31, where
nebular texture can correlate positively with seeing FWHM. The average global quality was
\num{0.469662}, ranging from \num{0.030536} to \num{0.820488}. The background-penalty input
averaged \num{25.636}, reflecting moderate sky-gradient variation across the frame sequence.

\begin{figure*}[t]
\centering
\includegraphics[width=\textwidth]{m31_aqmh_quality_timeseries.png}
\caption{Per-frame AQMH quality-map diagnostics. The shaded band shows the p10-p90 interval of
the map values; mean and median remain stable across the 645-frame sequence. The artifact
fraction stays near zero, with a higher early-frame tail.}
\label{fig:timeseries}
\end{figure*}

\begin{figure*}[t]
\centering
\includegraphics[width=\textwidth]{m31_aqmh_diagnostic_distributions.png}
\caption{Distributions of AQMH per-frame diagnostic quantities. The narrow quality-map mean
distribution reflects the within-frame relative normalization of the sigmoid factor; the
artifact gate remains an absolute veto term.}
\label{fig:distributions}
\end{figure*}

\subsection{Registration-weight guard}

The registration-weight guard was enabled and active. Of the 645 frames, 638 were registered
via \texttt{direct\_global}, 6 via \texttt{sequential\_refined}, and 1 was the reference frame.
The guard damped 6 frames (the sequential and reference frames), with a mean factor of
\num{0.9982} and a minimum factor of \num{0.7262}. The median factor was \num{1.0}, confirming
that the vast majority of frames were unaffected.

\subsection{Validation metrics}

The validation artifact reports an AQMH method run. The median input seeing estimate was
\num{12.329} px, while the median output FWHM was \num{5.268} px. This corresponds to a
\num{57.273}\% reduction of the measured FWHM:
\[
  100 \left(1 - \frac{5.268}{12.329}\right) = 57.273 .
\]
Background RMS changed from \num{4.185} to \num{0.780}, an \num{81.351}\% decrease.

The uniform-control gate compared the raw AQMH output against the unweighted mean. The results
are summarised in Table~\ref{tab:ucg}. The FWHM, seam score, elongation, and tail$_{11}$
regressions all passed their thresholds. The background-RMS regression (\num{0.168}) exceeded
the configured threshold of \num{0.05}, triggering the gate. As a result, the raw AQMH output
was preserved by the guard (\texttt{raw\_aqmh\_preserved\_by\_guard = true},
\texttt{selected\_candidate = raw\_aqmh}).

The low-frequency neutralisation candidate was evaluated but not applied: its background-RMS
regression (\num{0.168}) was not strictly smaller than the raw AQMH regression. The
structure-masked detail candidate was also evaluated but not applied: while its background-RMS
regression (\num{0.016}) was smaller, it did not pass all applicable validation thresholds
simultaneously. The alpha-blended attenuation search found no $\alpha > 0$ that satisfied all
gates, so the base (raw AQMH) was retained.

\begin{table}[t]
\centering
\caption{Uniform-control gate validation results.}
{\small
\begin{tabular}{@{}lrrr@{}}
\toprule
Metric & AQMH & Control & Regression \\
\midrule
FWHM (px) & 2.600 & 2.573 & $+0.010$ \\
Background RMS & 0.850 & 0.728 & $+0.168$ \\
Seam score & 1.736 & 1.830 & $-0.051$ \\
Elongation & 1.098 & 1.100 & $-0.001$ \\
Tail$_{11}$ abs & 0.492 & 0.494 & $-0.006$ \\
\bottomrule
\end{tabular}
}
\label{tab:ucg}
\end{table}

\begin{center}
\centering
\includegraphics[width=\columnwidth]{m31_aqmh_validation_summary.png}
\captionof{figure}{Validation summary for the M31 AQMH run. FWHM improves while the
background-RMS regression exceeds the gate threshold, causing the uniform-control gate to
trigger. The raw AQMH output is preserved.}
\label{fig:validation}
\end{center}

\begin{center}
\centering
\includegraphics[width=\columnwidth]{m31_aqmh_hms_preview.png}
\captionof{figure}{Final RGB preview after AQMH reconstruction, PCC, and HMS. The preview is a
visual run product derived from the FITS output, not a quantitative measurement.}
\label{fig:preview}
\end{center}

\section{Discussion}

\subsection{What the run demonstrates}

The M31 run demonstrates the intended AQMH v0.2.0 execution path with all post-reconstruction
extensions active:
\begin{itemize}
  \setlength{\itemsep}{0.05em}
  \setlength{\parsep}{0pt}
  \setlength{\parskip}{0pt}
  \setlength{\topsep}{0.15em}
  \item every usable frame produced a quality map;
  \item the PSF-FWHM sharpness proxy was used for global quality, avoiding the
        Laplacian-variance inversion on extended-emission targets;
  \item the registration-weight guard damped weakly-registered frames without affecting the
        majority;
  \item reconstruction consumed AQMH maps without classic tile-weight fallback;
  \item no quality map samples were missing during reconstruction;
  \item the uniform-control gate correctly identified the background-RMS regression and
        preserved the raw AQMH output;
  \item low-frequency neutralisation and structure-masked detail blending were evaluated but
        rejected by the validation gates, confirming that the guard mechanism works as designed;
  \item the measured FWHM and background RMS validation metrics improved.
\end{itemize}
The support counters are particularly important for implementation validation. A successful
AQMH-labeled run would be ambiguous if reconstruction silently fell back to missing-map behavior,
unweighted means, or classic tile weights. The combination of 645 written maps, zero missing map
samples, and \texttt{fallback\_to\_classic=false} makes the measured output attributable to the
AQMH path rather than to a mixed or compatibility path.

\subsection{Interpretation of the background-RMS gate}

The uniform-control gate triggered on the background-RMS regression of \num{0.168}, well above
the configured threshold of \num{0.05}. This means the AQMH-weighted output has a \num{16.8}\%
higher robust background RMS than the unweighted mean. This is not necessarily a defect: the
AQMH weighting preserves high-frequency structure that the unweighted mean smooths out, which can
increase the pixel-difference-based RMS estimator. The gate's purpose is to detect cases where
this increase is excessive; the threshold of \num{0.05} is conservative. The raw AQMH output was
preserved because no post-processed candidate improved all applicable metrics simultaneously.

\subsection{Interpretation of stable map means}

The frame-to-frame mean of $Q^{\mathrm{map}}$ is deliberately stable. The sharpness and SNR terms
are z-normalized inside each frame, so their sigmoid contribution is mostly a within-frame
relative field. Cross-frame quality discrimination comes from the global frame weight $G_{f,c}$
and from absolute veto terms such as the artifact gate. Therefore, a narrow distribution of map
means is not a defect; it is a consequence of separating within-frame spatial quality from
between-frame global quality.

\subsection{Acceleration}

The AQMH\_MAPS phase ran on the CPU with 8 parallel workers, while the AQMH\_RECONSTRUCTION phase
used a native CUDA backend on an NVIDIA GeForce GTX 1660 Ti (\SI{6}{GB} VRAM). The reconstruction
processed \num{23} chunks of \num{101} rows each, with a device memory budget of
\SI{2.33}{GB}. The total data volume was approximately \SI{52.8}{GB}, requiring
\num{4} sigma-clip iterations per chunk. The effective throughput was \SI{26.6}{MB/s}, indicating
a memory-bandwidth-limited pipeline on this GPU class.

\subsection{Limitations}

This paper reports one M31 run and should not be read as a full statistical validation across
object classes, sensors, exposure lengths, sky brightness levels, and seeing regimes. The
diagnostic flag \texttt{scene\_dependent\_snr=true} appears for all frames and deserves further
evaluation: the fallback is deterministic, but future work should quantify when
background-centered local SNR support is sufficient without the degraded-path flag. The
background-RMS gate triggered, which means the raw AQMH output has measurably higher
pixel-difference RMS than the uniform control; while this is expected for a quality-weighted
reconstruction that preserves high-frequency detail, the threshold calibration deserves
sensitivity analysis. The validation metrics are pipeline diagnostics rather than an externally
calibrated astronomical quality assessment. The final preview is not a scientific photometry
product; it includes downstream PCC and HMS processing.

Further validation should include repeated runs, independent datasets, explicit comparisons
against Classic Tile Compile and conventional stackers, sensitivity tests for pyramid parameters
and map resolution, and visual/quantitative checks around strong gradients, saturated stars,
nebular structure, and low-overlap frame borders.

\section{Conclusion}

AQMH v0.2.0 provides a deterministic pixel-wise quality-weighted alternative to coarse local tile
weighting for short-exposure deep-sky reconstruction. The M31 run confirms that the implementation
can execute the complete AQMH path on a 645-frame dataset, write all dense quality maps, apply the
registration-weight guard, evaluate the post-reconstruction extension suite (neutralisation,
structure-masked blending), and produce improved run-level validation metrics. The observed FWHM
reduction of \num{57.273}\%, the background RMS reduction of \num{81.351}\%, the absence of
missing map samples, and the correct operation of the uniform-control gate support the practical
viability of AQMH as an independent reconstruction method in tile\_compile.

The post-reconstruction extensions (background penalty, registration-weight guard, low-frequency
neutralisation, structure-masked detail blending) were all active and correctly evaluated. The
validation gates rejected both the neutralised and the structure-masked candidates, preserving
the raw AQMH output. This demonstrates that the guard mechanism works as designed: AQMH will not
silently replace its output with a post-processed variant that fails the documented regression
thresholds.

\newpage
\section*{Code and Data Availability}

\textbf{Reference method.}
Classic Tile Compile is described by
\href{https://doi.org/10.5281/zenodo.19437794}{\emph{Tile-Based Quality Reconstruction for DSO},
Version 3.3.9}. The corresponding repository source is the
\href{https://github.com/jeamy/tile_compile/blob/master/docs/v3/zenodo-3.3.9/paper-tile_based_quality_reconstruction_methodology_v_3.3.9_en.tex}{v3.3.9
paper source}.

\textbf{Previous version.}
The AQMH v0.1.0 paper bundle is archived in
\href{https://github.com/jeamy/tile_compile/tree/master/docs/AQMH/zenodo-0.1.0}{docs/AQMH/zenodo-0.1.0}.

\textbf{Methodology.}
The AQMH v0.2.1 methodology document is
\href{https://github.com/jeamy/tile_compile/blob/master/docs/AQMH/aqmh_methodik_en_v0.2.1.md}{aqmh\_methodik\_en\_v0.2.1.md}.

\textbf{Code.}
The public project repository is
\href{https://github.com/jeamy/tile_compile}{github.com/jeamy/tile\_compile}.

\textbf{This paper.}
The AQMH v0.2.0 paper bundle is archived in
\href{https://github.com/jeamy/tile_compile/tree/master/docs/AQMH/zenodo-0.2.0}{docs/AQMH/zenodo-0.2.0}.
The compact values used for the tables and plots are stored in
\href{https://github.com/jeamy/tile_compile/blob/master/docs/AQMH/zenodo-0.2.0/m31_aqmh_run_summary.json}{m31\_aqmh\_run\_summary.json}.

\textbf{Run data.}
The full source measurement artifacts are local outputs from an M31 analysis run and are not
part of the public repository bundle.

\end{document}
